% Standalone mathematical fragment. Requires amsmath, amssymb, amsthm.
% Independent foundational audit, 22 October 2026.
\section{The tensor foundation and the finite-lift implication}
\label{sec:tensor-audit}

The arguments in this section are complete and do not use the local C1/C2
certificates. They identify exactly what those certificates must supply.
The base field throughout is the full field $\mathbb R$.

\subsection{Two labelled fields and their residues}
Let $K_c,K_t$ be real closed extensions of $\mathbb R$ and put
\[
 T=K_c\otimes_{\mathbb R}K_t,\qquad
 B_i=\{x\in K_i: |x|<M\text{ for some }M\in\mathbb R_{>0}\},
 \quad i\in\{c,t\}.
\]
Write $\mathfrak m_i$ for the elements whose absolute value is smaller
than every positive real number. The order on $K_i$ is its unique field
order. No order, embedding into a field, or algebraic-disjointness
assumption is imposed on $T$.

Every $x\in B_i$ has a unique real standard part $\bar x$ with
$x-\bar x\in\mathfrak m_i$. Indeed, the real cut
$\{r\in\mathbb R:r\le x\}$ is nonempty and bounded; its supremum
has the claimed property. Two such real numbers differ by a real
infinitesimal and hence agree. Infinitesimals form an ideal of $B_i$,
so $x\mapsto\bar x$ is an $\mathbb R$-algebra homomorphism with
kernel $\mathfrak m_i$. In particular
$B_i=\mathbb R\oplus\mathfrak m_i$ as real vector spaces.

Since tensoring real vector spaces preserves injections,
$B:=B_c\otimes_{\mathbb R}B_t$ embeds into $T$. Its residue map is
\[
 \rho:B\longrightarrow\mathbb R,\qquad
 \rho\Bigl(\sum_i a_i\otimes b_i\Bigr)=\sum_i\bar a_i\bar b_i.
 \tag{TF1}
\]
The direct-sum decompositions above show that
\[
 \ker\rho=\mathfrak m_c\otimes B_t+B_c\otimes\mathfrak m_t.
 \tag{TF2}
\]
Thus every tensor of zero residue has a finite expression in which
each elementary term has at least one infinitesimal factor. None of
these observations requires the two field extensions to be disjoint.

\subsection{A strict real reserve gives an exact tensor sum of squares}
\paragraph{Lemma TF1 (positive residue).}
If $a\in B$ and $\rho(a)>0$, then $a$ is a finite sum of squares in
$T$. In fact its square roots may all be chosen to be elementary
tensors.

\begin{proof}
By (TF2), write
\[
 a=\varepsilon+\sum_{j=1}^N x_j\otimes y_j,\qquad
 \varepsilon=\rho(a)>0,
 \tag{TF3}
\]
with bounded pure factors and an infinitesimal in at least one factor
of each term. If $N=0$, the assertion is immediate. Otherwise choose
positive real numbers $\varepsilon_j$ summing to $\varepsilon$.
For a term with $x_j$ infinitesimal, choose a positive real
$M_j>|y_j|$. Direct multiplication gives
\[
 \varepsilon_j+x_j\otimes y_j
 =\frac12\left[
  (\varepsilon_j/M_j+x_j)\otimes(M_j+y_j)
 +(\varepsilon_j/M_j-x_j)\otimes(M_j-y_j)\right].
 \tag{TF4}
\]
All four pure factors are strictly positive in their respective
fields. Each elementary summand therefore is a square: for example
$\tfrac12u\otimes v=(\sqrt{u/2}\otimes\sqrt v)^2$.
If instead $y_j$ is infinitesimal, interchange the roles of the
two factors in this calculation. Sum (TF4) over $j$. Every operation
is a field operation in one labelled factor; no mixed radical or
mixed inverse occurs.
\end{proof}

\paragraph{Lemma TF2 (strict bounded Gram).}
Let $Q\in\operatorname{Sym}_r(B)$ and suppose its entrywise real
residue $\rho(Q)$ is positive definite. Then
\[
                 g^{\mathsf T}Qg\in\Sigma T^2
                 \qquad\text{for every }g\in T^r.
 \tag{TF5}
\]
The entries of $g$ need not be bounded, independent, or polynomial.

\begin{proof}
Choose a real invertible matrix $C$ with
$C^{\mathsf T}\rho(Q)C=I$. Replacing $Q$ by
$C^{\mathsf T}QC$ and $g$ by $C^{-1}g$ preserves the quadratic
expression, so assume $\rho(Q)=I$. Choose a real $\delta>0$ such
that $(r-1)\delta<1$. With $e_i$ the real coordinate vectors,
\[
\begin{split}
 Q={}&\sum_i\bigl(Q_{ii}-(r-1)\delta\bigr)e_ie_i^{\mathsf T}\\
 &+\sum_{i<j}\frac{\delta+Q_{ij}}2
                (e_i+e_j)(e_i+e_j)^{\mathsf T}\\
 &+\sum_{i<j}\frac{\delta-Q_{ij}}2
                (e_i-e_j)(e_i-e_j)^{\mathsf T}.
\end{split}
 \tag{TF6}
\]
The off-diagonal entries agree because the last two contributions
have difference $Q_{ij}$; their diagonal contributions sum to
$\delta$ for every incident pair. The residues of all scalar
coefficients in (TF6) are positive. Lemma TF1 expresses each
coefficient as a finite sum of squares in $T$. Evaluation on $g$
then expresses (TF5) as a finite sum of squares of those scalar roots
times $g_i$ or $g_i\pm g_j$. This is a polynomial identity in the
commutative ring $T$, so no size condition on $g$ is involved.
\end{proof}

\paragraph{The strictness cannot be dropped.}
Let $K$ be a real closure of $\mathbb R(\tau)$ ordered so that
$\tau>0$ is infinitesimal. The order-preserving automorphism
$\tau\mapsto2\tau$ of $\mathbb R(\tau)$ extends to an
order-preserving automorphism $\sigma$ of $K$. In $K\otimes_{\mathbb R}K$,
\[
                  D=\tau\otimes1-1\otimes\tau
 \tag{TF7}
\]
has zero residue and zero value under diagonal multiplication.
Nevertheless the ring homomorphism
$x\otimes y\mapsto x\sigma(y)$ sends it to $-\tau<0$.
Consequently $D$ is not a sum of squares. This example also shows
why positivity of a source and test after identifying the two tensor
factors cannot replace an identity in $T$. It is a boundary example
for the lemma, not an asserted counterexample to a Bell slack.

\subsection{Legal normalization and finite assembly}
\paragraph{Lemma TF3 (pure-factor cancellation).}
For $s\in K_c^{>0}$, $t\in K_t^{>0}$, and $F\in T$,
\[
 F\in\Sigma T^2\quad\Longleftrightarrow\quad
                  (s\otimes t)F\in\Sigma T^2.
 \tag{TF8}
\]
Indeed multiply each root by $\sqrt s\otimes\sqrt t$ in one
direction and by $s^{-1/2}\otimes t^{-1/2}$ in the other.
These are inverse units of $T$. Nonvanishing must be proved before
this lemma is used. A homogeneous coordinate that is literally zero
cannot be cancelled, and the lemma does not authorize cancellation
of a general mixed tensor.

\paragraph{Lemma TF4 (positive matrix pairing).}
If $X\in\operatorname{Sym}_r(K_c)$ and
$Y\in\operatorname{Sym}_r(K_t)$ are positive semidefinite in their
own fields, then
\[
 \sum_{i,j}X_{ij}\otimes Y_{ij}
 =\sum_{a,b}\left(\sum_i U_{ia}\otimes V_{ib}\right)^2,
 \qquad X=UU^{\mathsf T},\quad Y=VV^{\mathsf T}.
 \tag{TF9}
\]
To obtain the factorizations over a real closed field, pivot on a
positive diagonal entry and complete its square. The remaining
Schur complement is positive semidefinite. If every diagonal entry
is zero, every off-diagonal entry is zero by its two-dimensional
quadratic restriction. Induction, and positive square roots of the
pivots, finish the factorization even for singular matrices.
Expansion of the right side proves (TF9).

\paragraph{Lemma TF5 (conic and convex closure).}
For $f=f_0+\sum_{\nu=1}^n f_\nu q_\nu$ set
$[f,q]=f_0\otimes1+\sum_\nu f_\nu\otimes q_\nu$.
Suppose in the source field
$f=\sum_j a_j f_j$ with $a_j\ge0$, and in the testing field
$q=\sum_i w_i q_i$ with $w_i\ge0$ and $\sum_iw_i=1$.
If $[f_j,q_i]=\sum_k h_{ijk}^2$ in $T$, then
\[
             [f,q]=\sum_{i,j,k}
                \left[(\sqrt{a_j}\otimes\sqrt{w_i})h_{ijk}\right]^2.
 \tag{TF10}
\]
This follows from linearity in $f$ and affineness in $q$; the equation
$\sum_iw_i=1$ preserves the constant coefficient. All sums are finite.
Zero weights give zero roots and are never inverted.

\subsection{Exactly the external criterion that is needed}
Let $S\subset\mathbb R^n$ be closed, convex, and semialgebraic, and
let $P_S$ be its cone of nonnegative affine functions. Fix a finite
real formula defining $S$. Over any real closed extension $K$, its
transfer $S_K$ is the set defined by that same formula. The formula
\[
          \forall q\ (q\in S\ \Longrightarrow\ f(q)\ge0)
 \tag{TF11}
\]
shows, by quantifier elimination and elementary transfer, that
$(P_S)_K$ consists exactly of the $K$-affine functions nonnegative
on $S_K$. This is a statement about formulas; it makes no claim
that $S_K$ is compact in the order topology of $K$.

Scheiderer's Corollary 3.19 says that $S$ has a finite semidefinite
lift if and only if, for every real closed extension $K/\mathbb R$,
every $f\in(P_S)_K$, and every $q\in S_K$, the point-first tensor
evaluation is a finite sum of squares in $K\otimes_{\mathbb R}K$.
The coefficient-first convention used here is equivalent by the
factor-flip automorphism, which preserves squares.
Therefore a proof of
\[
       [f,q]\in\Sigma(K\otimes_{\mathbb R}K)^2
       \quad\text{for every }K,\ f\in(P_S)_K,\ q\in S_K
 \tag{TF12}
\]
is sufficient. Two labelled copies of each one arbitrary $K$ suffice;
a theorem for all unequal pairs $K_c,K_t$ is stronger than necessary.
The labels must still remain separate in every identity. In particular,
algebraic disjointness of the copies must not be imposed.

The quantifiers in (TF12) permit the number and tensor ranks of roots
to depend on $K,f,q$. Corollary 3.19 supplies finiteness of a common
lift through its finite-dimensional algebraic construction and
compactness argument. No uniform root-count or rank bound has to be
supplied beforehand. Theorem 3.10 is the stronger quantitative version:
an asserted block-order bound $d$ needs a uniform bound $d$ on the
tensor ranks of the square roots. The lemmas above do not by themselves
provide such a bound for all local certificates. Both results concern
arbitrary finite ambient dimension, despite the paper's planar title.

\paragraph{Primary source and exact locations.}
C.~Scheiderer, \emph{Second-order cone representation for convex sets
in the plane}, arXiv:2004.04196v2 (28 January 2021): the bounded-ring
setup is Section 3.2, printed p.~10; Proposition 3.5 is on printed
p.~11; Theorem 3.10 is on printed p.~12; Corollary 3.19 and its proof
are on printed p.~15. The elementary reserve identity (TF4) above
independently proves the positive-residue step for the two-labelled-field
setting. The external implication used is Corollary 3.19 itself.

\paragraph{Audit conclusion.}
The foundational tensor arguments in the recovered October manuscript
are valid under their stated strictness and field hypotheses. They
reduce the remaining verification to actual full-parent local
identities, boundedness and positive-definite residue of their stated
Grams, and complete source/test coverage. These lemmas do not establish
any of those remaining domain-specific assertions by themselves.

\subsection{What a determinant certificate does and does not supply}
Here is a useful lawful reduction, together with its sharp limitation.
If $P\in\operatorname{Sym}_4(B)$ has positive semidefinite rank-three
residue, put
\[
 a=e_3(P),\qquad S=P^2-e_1(P)P+e_2(P)I.
\]
Cayley--Hamilton gives the exact identity
\[
 a\,x^TPx=(Px)^TS(Px)+\det(P)\,x^Tx.                 \tag{TF13}
\]
If the nonzero residue eigenvalues are $\lambda_1,\lambda_2,\lambda_3$,
then $\rho(a)=\lambda_1\lambda_2\lambda_3>0$, and the four
eigenvalues of $\rho(S)$ are
$\lambda_2\lambda_3,\lambda_1\lambda_3,\lambda_1\lambda_2,
\lambda_1\lambda_2+\lambda_1\lambda_3+\lambda_2\lambda_3$.
Thus a determinant decomposition $\det(P)=\sum_\nu q_\nu^2$
gives a bounded strict Gram for $aF$ in the frame
$z=((Px)_i,(q_\nu x_j)_{\nu,j})$, where $F=x^TPx$.

\paragraph{Lemma TF6 (bounded conductor cancellation).}
Suppose $aF=z^TGz$ with $a\in B$, $\rho(a)>0$, and
$G\in\operatorname{Sym}_r(B)$ of positive definite residue.
If also $F=z^THz$ for some $H\in\operatorname{Sym}_r(B)$, then
$F$ is tensor SOS. Indeed, with $a_0=\rho(a)>0$,
\[
 a_0F=z^T\{G-(a-a_0)H\}z.                           \tag{TF14}
\]
The displayed Gram has residue $\rho(G)$, so Lemma TF2 applies.
Only the positive real number $a_0$ is divided out. This reduces
positivity to a bounded quadratic membership, but does not establish
that membership.

\paragraph{Counterexample to automatic bounded membership.}
The following recovered R17 example is independently checked here.
Let $K_s$ be a real closure of $\mathbb R(s)$ with $s>0$
infinitesimal, let $K_c$ be the real closure of $K_s(\varepsilon)$
with $0<\varepsilon<r$ for every $r>0$ in $K_s$, and let $K_t$
be a real closure of $\mathbb R(t,v)$ with $t,v>0$ infinitesimal.
Write separated factors by juxtaposition, and set
\[
 A=1+s^2t^2,\qquad
 c=\varepsilon^2\{(t-s)^2+(v-s)^2\},\qquad
 P=\operatorname{diag}(A,1,1,c).                     \tag{TF15}
\]
Both $x^TPx$ and $\det(P)=Ac$ are explicit finite tensor SOS.
The residue matrix is $\operatorname{diag}(1,1,1,0)$.
In scalar source-field variables there is an isolated stationary
zero at $(t,v,x)=(s,s,e_4)$, with positive angular curvatures
$2\varepsilon^2$ and positive three-dimensional state normal form.

Nevertheless, at the independent state $x=e_1$, no determinant-root
spectral frame $z=((Px)_i,(q_\nu x_j))$ admits a bounded identity
$F=z^THz$, for \emph{any} finite SOS choice $\det(P)=\sum q_\nu^2$.
To prove this, use the coarser $\varepsilon$-valuation, whose residue
field on $K_c$ is $K_s$. Every coefficient bounded by a real number
has nonnegative valuation; its valuation-zero residue is bounded
in $K_s$. Embed $K_c$ in the real Puiseux field over $K_s$.
Finite tensor sums then have Puiseux leading coefficients in
$T_0=K_s\otimes_{\mathbb R}K_t$, a formally real domain.
Leading squares cannot cancel, so valuation $2$ of $\det(P)$
forces every $q_\nu$ to have valuation at least $1$.
Taking valuation-zero terms in the proposed bounded identity gives
\[
                  A=A^2h,\qquad h\in B_s\otimes B_t. \tag{TF16}
\]
But $A$ is not a unit even in $T_0$: in an algebraically closed
overfield containing $K_s$, send $t$ to $i/s$ and send $v$ to a
new independent transcendental. This embeds $\mathbb R(t,v)$
and extends to its algebraic real closure $K_t$. The resulting
homomorphism on $T_0$ sends $A$ to zero and $1$ to $1$.
Since $T_0$ is a domain and $A\ne0$, (TF16) is impossible.

This example already has an SOS parent. It refutes automatic
bounded-conductor extraction from PSD, determinant SOS, strict
normal rank, stationary contact and rank-one positive event
decompositions. It does \emph{not} refute tensor SOS of the Q222
parent, or every conceivable matrix lifting theorem. A minimal
proof using (TF13)--(TF14) must still supply a source-specific
bounded membership, possibly in an enlarged event-adapted frame.
In particular, even the mixed tensor $1+s^2t^2$, a sum of two
positive elementary tensors with residue one, is not invertible.
Calling such a pivot positive never licenses mixed division.
